% ===========================================================
%  第 8 章 相似理论 —— 北大《高等代数》第五版【深化】补充
%  源：北大《高等代数》第五版（王萼芳、石生明）
%  页码换算：PDF 页 = 印刷页 + 12（印刷 1–163 区间已实测）
%  本片与 deep_ch08_authored.tex（主控手写 3 块）互不重复：
%    authored 已覆盖 相似=换基 / 可对角化五条件 / Jordan 块为什么不能对角化
%    本片补充：Cayley–Hamilton、最小多项式、不变子空间、相似不变量清单、
%              实对称矩阵正交对角化、正交变换化二次型
% ===========================================================

% ---------- 精读页面：p215–p218(§7.5) p221–p224(§7.7) p225–p227(§7.8)
%              p228–p230 + PDF231(§7.9 收尾) p269–p274(§9.6) ----------

% ==================== §7.5 对角矩阵（线索起点） ====================

\begin{fdeep}{哈密顿--凯莱定理：最小多项式为什么存在}
设 $A$ 是 $n$ 阶数字矩阵，$f(\lambda)=|\lambda E-A|=\lambda^{n}+a_{1}\lambda^{n-1}+\dots+a_{n-1}\lambda+a_{n}$，则
\fml{f(A)=A^{n}+a_{1}A^{n-1}+\dots+a_{n-1}A+a_{n}E=O}
即\textbf{每个矩阵都满足自己的特征多项式}。对线性变换 $\mathscr{A}$ 同理：$f(\mathscr{A})=\mathscr{O}$。

\textbf{证明思路（易错点所在）}：设 $B(\lambda)=\lambda^{n-1}B_{0}+\lambda^{n-2}B_{1}+\dots+B_{n-1}$ 为
$\lambda E-A$ 的伴随矩阵，则 $B(\lambda)(\lambda E-A)=f(\lambda)E$。比较 $\lambda$ 同次幂的系数得到
$n+1$ 个矩阵等式，再用 $A^{n},A^{n-1},\dots,A,E$ \textbf{依次从右边乘}第 $1,2,\dots,n+1$ 式，
把 $n+1$ 式\textbf{相加}，左边两两抵消为 $O$，右边正好是 $f(A)$。
\fml{\text{不能直接把 }\lambda=A\text{ 代入}}
因为 $B(\lambda)$ 的元素是 $\lambda$ 的多项式，代入后 $B(\lambda)$ 无法定义。

【反哺点】服务考纲"理解…可相似对角化的充分必要条件"。$f(A)=O$ 保证"以 $A$ 为根的非零多项式"
至少有一个，最小多项式才有定义；同时立刻得到 $m(\lambda)\mid f(\lambda)$（$m$ 指最小多项式）。
考场上求 $A^{100}$ 一类题，最省力的解法也是拿 $f(A)=O$ 降次，而非硬算。
\end{fdeep}

\begin{fdeep}{最小多项式的定义与四条基本性质}
\textbf{定义}：数域 $P$ 上使 $f(A)=O$ 的多项式称为以 $A$ 为根的多项式，其中\textbf{次数最低、
首项系数为 $1$} 的那一个称为 $A$ 的\textbf{最小多项式}，记作 $m(\lambda)$。
（由哈密顿--凯莱定理，这样的多项式一定存在。）

\textbf{四条基本性质}：
\begin{enumerate}[leftmargin=2.2em,itemsep=0.2em]
\item \textbf{唯一性}：$m(\lambda)$ 唯一。（若 $g_1,g_2$ 都是，由带余除法 $g_1=qg_2+r$ 与
      $r(A)=O$、$\partial r<\partial g_2$ 得 $r=0$。）
\item \textbf{整除判据}：$f(A)=O\iff m(\lambda)\mid f(\lambda)$，
      这是\textbf{计算最小多项式最好用的一条}，只需验证候选多项式是否以 $A$ 为根。
\item \textbf{整除特征多项式}：$m(\lambda)$ 是特征多项式 $f(\lambda)$ 的因式。
\item \textbf{相似不变量}：若 $B=T^{-1}AT$，则 $f(B)=T^{-1}f(A)T$，故 $f(B)=O\iff f(A)=O$，
      \textbf{相似矩阵的最小多项式相同}。
\end{enumerate}
\textbf{高频特例}：$kE$ 的是 $\lambda-k$，单位矩阵的是 $\lambda-1$，零矩阵的是 $\lambda$；
反之\textbf{最小多项式为 1 次多项式 $\Rightarrow$ $A$ 是数量矩阵}。

【反哺点】服务考纲"理解…可相似对角化的充分必要条件"。性质 2 是求 $m(\lambda)$ 的\textbf{操作准则}：
先分解 $f(\lambda)$，再按次数从低到高逐个试 $f(A)=O$。"相似矩阵最小多项式相同"又给出一个
相似不变量，可直接用于"是否相似"的判断。
\end{fdeep}

\begin{fdeep}{最小多项式的计算规则：与准对角、Jordan 块的关系}
\textbf{规则一（准对角）}：设 $A=\operatorname{diag}(A_1,A_2,\dots,A_s)$，$A_i$ 的最小多项式为 $g_i(\lambda)$，则
\fml{m_A(\lambda)=[\,g_1(\lambda),\,g_2(\lambda),\,\dots,\,g_s(\lambda)\,]}
即各块最小多项式的\textbf{最小公倍式}。证明要点：$g=[g_1,\dots,g_s]$ 使 $g(A)=O$；反之若
$h(A)=O$，则 $h(A_i)=O$，故每个 $g_i\mid h$，于是 $g\mid h$。

\textbf{规则二（$k$ 阶若尔当块）}：$J_k(a)$ 的最小多项式是 $m(\lambda)=(\lambda-a)^{k}$；
其特征多项式同样是 $(\lambda-a)^k$，但 $(J_k(a)-aE)^{k-1}\ne O$，故次数不能再降。
由此得一\u200b条极好用的计数规律：$(\lambda-\lambda_0)$ 在 $m(\lambda)$ 中的重数
$=$ 特征值 $\lambda_0$ 对应的\textbf{最大 Jordan 块阶数}；而在 $f(\lambda)$ 中的重数
$=$ 该特征值全部 Jordan 块阶数之和（代数重数）。

【反哺点】服务考纲"理解…矩阵可相似对角化的充分必要条件"。规则一把"分块矩阵求最小多项式"
变成一步最小公倍式；规则二给出的"最小多项式重数 = 最大块阶数"是判定可对角化、
以及判定两个矩阵是否相似时最锋利的工具。
\end{fdeep}

\begin{fdeep}{最小多项式无重根 $\iff$ 可相似对角化}
\textbf{定理（北大 §7.9 定理 15 及推论）}：数域 $P$ 上 $n$ 阶矩阵 $A$ 与对角矩阵相似的
\textbf{充分必要条件}是 $A$ 的最小多项式是 $P$ 上互素的一次因式的乘积；在复数域上即
\fml{A\sim\Lambda\iff m(\lambda)\text{ 无重根}}
显式写法为 $m(\lambda)=(\lambda-\lambda_1)(\lambda-\lambda_2)\dots(\lambda-\lambda_s)$，
其中 $\lambda_1,\dots,\lambda_s$ 是 $A$ 的全部\textbf{互不相同}的特征值。

\textbf{两半的证明思路}：\textbf{必要性}——对角阵的最小多项式是各对角元 $\lambda_i$ 对应的一次
因式 $(\lambda-\lambda_i)$ 的最小公倍式，互不相同的一次因式其最小公倍式即为乘积，故无重根。
\textbf{充分性}——设 $g(\lambda)=\prod_{i=1}^{s}(\lambda-a_i)$，由 $g(\mathscr{A})=\mathscr{O}$ 知
$V=\ker g(\mathscr{A})$；把 $g$ 的 $s$ 个\textbf{互素}一次因式对应的核作直和分解
$V=V_1\oplus\dots\oplus V_s$，每个基向量都落在某个 $V_i$ 中因而是特征向量。
关键技术是\textbf{互素}保证贝祖等式 $\sum u_i(\lambda)f_i(\lambda)=1$，
从而每个向量都能拆成各核中向量之和；而一次因式使 $V_i$ 恰好是特征子空间。

【反哺点】服务考纲"理解…矩阵可相似对角化的充分必要条件"。这是所有判据中\textbf{唯一既不需要
解特征向量、也不需要数几何重数}的一条：含参数矩阵只要算出 $m(\lambda)$ 看有无重根即可。
反过来"最小多项式有重根"就是\textbf{不可对角化的判决书}，比逐个特征值算秩更省时间。
\end{fdeep}

% ==================== §7.7 不变子空间 ====================

\begin{fdeep}{不变子空间：可对角化的几何解释}
\textbf{定义}：设 $\mathscr{A}$ 是数域 $P$ 上线性空间 $V$ 的线性变换，$W\le V$。若对 $W$ 中
任一向量 $\xi$ 都有 $\mathscr{A}\xi\in W$，则称 $W$ 是 $\mathscr{A}$ 的\textbf{不变子空间}。

\textbf{四类基本例子}：整个空间 $V$ 与零子空间 $\{0\}$ 对任何 $\mathscr{A}$ 都不变；
$\mathscr{A}$ 的\textbf{值域 $\mathscr{A}V$ 与核 $\mathscr{A}^{-1}(0)$} 都是不变子空间；
若 $\mathscr{A}$ 与 $\mathscr{B}$ 可交换，则 $\mathscr{B}$ 的核与值域都是 $\mathscr{A}$-子空间
——更一般地，\textbf{$\mathscr{A}$ 的任何多项式 $f(\mathscr{A})$ 的值域与核}都是 $\mathscr{A}$-子空间；
任何子空间都是\textbf{数乘变换}的不变子空间。

\textbf{一维不变子空间与特征向量完全是一回事}：设 $W=L(\xi)$ 是一维 $\mathscr{A}$-子空间，
$\xi\ne0$，则 $\mathscr{A}\xi\in W$ 即 $\mathscr{A}\xi=\lambda_0\xi$——$\xi$ 就是特征向量；
反之特征向量 $\xi$ 的倍数集合 $L(\xi)$ 是一维 $\mathscr{A}$-子空间。

【反哺点】服务考纲"理解…矩阵可相似对角化的充分必要条件"。本块把"可对角化"翻译成几何语言：
\textbf{$A$ 可对角化 $\iff$ $V$ 能分解为一维不变子空间的直和}。特征子空间 $V_{\lambda_0}$
作为整体也是不变子空间，这就是"特征方向"的严格说法——它解释了我们关心的为什么是
特征向量的\textbf{个数}，而不是特征值的个数。
\end{fdeep}

\begin{fdeep}{不变子空间与矩阵形状：分块三角、准对角}
\textbf{结论一（分块三角）}：设 $W$ 是 $n$ 维空间 $V$ 的 $\mathscr{A}$-子空间，$\dim W=k$。
取 $W$ 的基 $\varepsilon_1,\dots,\varepsilon_k$ 并扩充为 $V$ 的基 $\varepsilon_1,\dots,\varepsilon_n$，
则 $\mathscr{A}$ 的矩阵呈分块三角形状 $\begin{pmatrix}A_1&A_3\\O&A_2\end{pmatrix}$，且左上角 $k$ 阶块
$A_1$ 正是限制变换 $\mathscr{A}|_W$ 的矩阵；反之若矩阵呈此形状，则由 $\varepsilon_1,\dots,\varepsilon_k$
生成的子空间必是 $\mathscr{A}$-子空间。（限制变换与 $\mathscr{A}$ 的区别：$(\mathscr{A}|_W)\xi=\mathscr{A}\xi$
只对 $\xi\in W$ 有意义。）

\textbf{结论二（准对角）}：若 $V=W_1\oplus\dots\oplus W_s$，每个 $W_i$ 都不变，在每个 $W_i$ 中取基
并合并为 $V$ 的基，则 $\mathscr{A}$ 的矩阵为准对角形 $\operatorname{diag}(A_1,\dots,A_s)$，
其中 $A_i$ 是 $\mathscr{A}|_{W_i}$ 的矩阵；反之亦然。于是
\fml{\textbf{矩阵分解为准对角形}\iff\textbf{空间分解为不变子空间的直和}}

【反哺点】服务考纲"掌握将矩阵化为相似对角矩阵的方法"。结论二把"找可逆 $P$ 使 $P^{-1}AP$ 为
对角阵"翻译成"把空间拆成一维不变子空间"，于是"为什么取特征向量当新基"变得显然。
\end{fdeep}

\begin{fdeep}{根子空间分解：一般情形下能保证的只是准对角}
\textbf{定理 12}：设线性变换 $\mathscr{A}$ 的特征多项式在数域 $P$ 上可分解为一次因式之积
$f(\lambda)=(\lambda-\lambda_1)^{r_1}\dots(\lambda-\lambda_s)^{r_s}$，则 $V$ 可分解为不变子空间的直和
\fml{V=V_1\oplus V_2\oplus\dots\oplus V_s,\qquad V_i=\{\xi\in V\mid(\mathscr{A}-\lambda_i\mathscr{E})^{r_i}\xi=0\}}
称 $V_i$ 为 $\mathscr{A}$ 的属于 $\lambda_i$ 的\textbf{根子空间}（记 $V^{\lambda_i}$）。

\textbf{证明思路（贝祖等式）}：令 $f_i(\lambda)=f(\lambda)/(\lambda-\lambda_i)^{r_i}$，取 $V_i=f_i(\mathscr{A})V$。
因 $f_1,\dots,f_s$ \textbf{互素}，存在多项式 $u_1,\dots,u_s$ 使 $\sum u_if_i=1$，于是
$\alpha=\sum u_i(\mathscr{A})f_i(\mathscr{A})\alpha$，每项都落在对应 $V_i$ 中——这就是"每个向量都能拆开"。

\textbf{由此得到的图像}：一般情况下能保证的只是\textbf{准对角}
（这正是若尔当形的来源）；只有当每个根子空间都恰好等于特征子空间时才退化为对角形。

【反哺点】服务考纲"理解…矩阵可相似对角化的充分必要条件"。本块交代了 $g_\lambda=m_\lambda$ 的
几何来历：根子空间 $V_i$ 的维数就是\textbf{代数重数} $r_i$，而其中真正被固定住的方向只有
特征子空间那些（$g_{\lambda_i}$ 个）。两者相等时根子空间就"化"成了特征子空间——这就是可对角化的完整图像。
\end{fdeep}

% ==================== §7.8 若尔当标准形（结构性补充） ====================

\begin{fdeep}{若尔当标准形的存在与唯一性（北大 §7.8 定理 13、14）}
\textbf{定理 13（存在性）}：设 $\mathscr{A}$ 是复数域上 $n$ 维空间 $V$ 的线性变换，则 $V$ 中一定
存在一组基，使 $\mathscr{A}$ 在这组基下的矩阵是若尔当形矩阵，称为 $\mathscr{A}$ 的
\textbf{若尔当标准形}。

\textbf{证明的两步骨架}：先用根子空间分解 $V=V_1\oplus\dots\oplus V_s$ 化归到单个根子空间；
在 $V_i$ 上 $(\mathscr{A}-\lambda_i\mathscr{E})^{r_i}=\mathscr{O}$，即 $\mathscr{B}=\mathscr{A}-\lambda_i\mathscr{E}$
是\textbf{幂零变换}。对幂零变换可证：$V$ 中必有一组形如
$\alpha_1,\mathscr{B}\alpha_1,\dots,\mathscr{B}^{k_1-1}\alpha_1,\ \alpha_2,\mathscr{B}\alpha_2,\dots$
的\textbf{循环基}（每条链尾满足 $\mathscr{B}^{k_i}\alpha_i=0$），它使 $\mathscr{B}$ 的矩阵为若尔当形；
再叠加 $\lambda_i\mathscr{E}$ 即得结论。

\textbf{定理 14（唯一性）}：每个 $n$ 阶复矩阵 $A$ 一定与一个若尔当形矩阵相似；这个若尔当形矩阵
\textbf{除去其中若尔当块的排列顺序外由 $A$ 唯一决定}。即若尔当形是三角形矩阵，
主对角线上的元素就是特征多项式的全部根，且\textbf{两个复矩阵相似 $\iff$ 它们的若尔当标准形
只差块的排列顺序}。

【反哺点】服务考纲"理解相似矩阵的概念、性质"。本块回答"不可对角化时矩阵究竟长什么样"：
不是"没有标准形"，而是标准形里出现了阶数 $\ge2$ 的块。由此也说明相似判定的\textbf{完整}
不变量是若尔当形，而迹、行列式、特征多项式都只是它的"部分投影"。
\end{fdeep}

% ==================== 相似不变量：完整清单 ====================

\begin{fdeep}{相似不变量的完整清单（以及哪一条才够用）}
把前面各章零散提到的相似不变量汇总，考场按此顺序逐条排除：
\begin{center}
\begin{tabular}{@{}ll@{}}
\textbf{相似不变量} & \textbf{作用} \\
\hline
特征多项式 $|\lambda E-A|$（含特征值及重数） & 最基础的筛子；不同则必不相似 \\
行列式 $|A|$、迹 $\operatorname{tr}(A)$ & 分别是特征多项式的常数项与 $n-1$ 次项系数，是推论 \\
秩 $r(A)$ & 部分反映特征值信息，零特征值个数需另算 \\
最小多项式 $m(\lambda)$ & 能抓到"最大若尔当块阶数"的信息 \\
若尔当标准形（不计块序） & \textbf{完全不变量}：同 $\iff$ 相似 \\
\end{tabular}
\end{center}
\textbf{两个必须警惕的点}：行列式、迹、秩单独任意一条或几条相同，都\textbf{不足以}推出相似；
特征多项式相同\textbf{也不够}——例如三阶矩阵 $\operatorname{diag}(1,1,1)$ 与 $J_3(1)$ 的特征
多项式都是 $(\lambda-1)^3$，但不相似。判定的\textbf{终点}是比较若尔当标准形，
或同时比较特征多项式与最小多项式。

【反哺点】服务考纲"理解相似矩阵的概念、性质及矩阵可相似对角化的充分必要条件"。
选择题常设"$A,B$ 特征值相同 $\Rightarrow$ $A\sim B$"这类陷阱；本清单给出正确的取舍顺序：
先用\textbf{特征多项式 + 秩}快速排除，再用\textbf{最小多项式}区分同特征多项式的情形，
最后才动用若尔当形。另外"可对角化性是一个相似不变量"这一条，
让"A 与 B 相似而 A 可对角化"类判断题可以只看一方。
\end{fdeep}

% ==================== §9.6 实对称矩阵的标准形 ====================

\begin{fdeep}{实对称矩阵：特征值全实、特征向量正交、必可正交对角化}
对应考纲"掌握实对称矩阵的特征值和特征向量的性质"，北大 §9.6 把三条性质串成一条链：
\begin{enumerate}[leftmargin=2.2em,itemsep=0.2em]
\item \textbf{引理 1}：实对称矩阵 $A$ 的\textbf{复特征值皆为实数}；
\item \textbf{引理 4}：实对称矩阵\textbf{属于不同特征值的特征向量必正交}；
\item \textbf{定理 7}：对任意 $n$ 阶实对称矩阵 $A$，都存在 $n$ 阶\textbf{正交矩阵 $T$}，使
      $T^{\mathrm{T}}AT=T^{-1}AT$ 成对角形。
\end{enumerate}
链的逻辑：引理 1 保证特征值都是实数，才谈得上实数域上的特征向量；引理 4 保证
"不同特征值的特征子空间已经两两正交"，于是把它们拼起来自动得到正交基，
再在每个特征子空间内部正交化即可，\textbf{无需整体施密特正交化}。

【反哺点】服务考纲"掌握实对称矩阵的特征值和特征向量的性质"。这三条也是"实对称矩阵为什么
一定可对角化"的完整答案：它的特征子空间之间\textbf{天然正交}，不会出现 Jordan 链那种
"几何重数不够"的病态。考场上凡遇实对称矩阵可直接断言必可正交对角化，省掉全部判别步骤；
而"不同特征值的特征向量正交"是解实对称矩阵题时的\textbf{免费检查手段}。
\end{fdeep}

\begin{fdeep}{实对称矩阵三条性质的证明要点（每条一个恒等式）}
\textbf{引理 1（特征值全实）}：设实对称矩阵满足 $A\xi=\lambda_0\xi$（$\xi\ne0$），考察
\fml{\overline{\xi}^{\mathrm{T}}(A\xi)=\overline{\xi}^{\mathrm{T}}A^{\mathrm{T}}\xi=(A\overline{\xi})^{\mathrm{T}}\xi}
左边为 $\lambda_0\overline{\xi}^{\mathrm{T}}\xi$，右边为 $\overline{\lambda_0}\,\overline{\xi}^{\mathrm{T}}\xi$；
而 $\overline{\xi}^{\mathrm{T}}\xi=\bar x_1x_1+\dots+\bar x_nx_n\ne0$，故 $\lambda_0=\overline{\lambda_0}$ 是实数。

\textbf{引理 4（特征向量正交）}：设 $A\alpha=\lambda\alpha$，$A\beta=\mu\beta$，$\lambda\ne\mu$。
由 $A$ 的对称性 $\beta^{\mathrm{T}}(A\alpha)=\alpha^{\mathrm{T}}A\beta$ 得
$\lambda(\alpha,\beta)=\mu(\alpha,\beta)$，因 $\lambda\ne\mu$，故 $(\alpha,\beta)=0$。

\textbf{定理 7（必可正交对角化）}：对维数 $n$ 归纳。取 $A$ 的单位特征向量 $\alpha_1$，记
$V_1=L(\alpha_1)$。关键是\textbf{对称变换的不变子空间的正交补仍是不变子空间}：若
$\alpha\perp V_1$、$\beta\in V_1$，则 $(\mathscr{A}\alpha,\beta)=(\alpha,\mathscr{A}\beta)=0$，
故 $\mathscr{A}\alpha\in V_1^{\perp}$，于是 $\mathscr{A}|_{V_1^{\perp}}$ 仍对称，归纳即得。

【反哺点】服务考纲"掌握实对称矩阵的特征值和特征向量的性质"。三个证明各只需一个恒等式，
记住它们，考场上遇到"证明实对称矩阵某性质"的题直接套用即可。
\end{fdeep}

\begin{fdeep}{求正交矩阵 $T$ 使 $T^{\mathrm{T}}AT$ 成对角形的三步法}
北大 §9.6 的算法：
\begin{enumerate}[leftmargin=2.2em,itemsep=0.2em]
\item 求出 $A$ 的全部不同的特征值 $\lambda_1,\dots,\lambda_s$；
\item 对每个 $\lambda_i$ 解 $(\lambda_iE-A)x=0$，其基础解系是 $V_{\lambda_i}$ 的一组基；
      由此求出 $V_{\lambda_i}$ 的一组\textbf{标准正交基} $\eta_{i1},\dots,\eta_{id_i}$；
\item 引理 4 保证不同特征值的向量两两正交，定理 7 保证向量总个数等于 $n$，
      于是这些 $\eta$ 构成 $\mathbb{R}^n$ 的标准正交基，且都是 $A$ 的特征向量。
\end{enumerate}
把它们按列排成矩阵即得所求正交矩阵
\fml{T=(\eta_1,\dots,\eta_n),\qquad T^{-1}AT=T^{\mathrm{T}}AT=\operatorname{diag}(\lambda_1,\dots,\lambda_n)}
\textbf{两个实用补丁}：① \textbf{正交化只需在同一特征子空间内部做}——不同特征值的特征向量
已自动正交，把整组 $n$ 个向量从头做施密特正交化是纯粹的浪费。② \textbf{总可要求 $|T|=1$}：
若 $|T|=-1$，取 $S=\operatorname{diag}(-1,1,\dots,1)$，则 $T_1=TS$ 仍正交、$|T_1|=1$
且 $T_1^{\mathrm{T}}AT_1=T^{\mathrm{T}}AT$ 仍为对角形。

【反哺点】服务考纲"掌握将矩阵化为相似对角矩阵的方法"与"掌握实对称矩阵…的性质"。
\textbf{正交化只在每个特征子空间内部做}，能省掉大量根式运算——这是实对称矩阵正交对角化
与一般矩阵相似对角化在操作上的最大差别；$T^{-1}=T^{\mathrm{T}}$ 也是唯一能"免费求逆"的场景。
\end{fdeep}


% ===========================================================================
%  【主控裁决 · 第 8 章 deep_ch08_beida 取舍】（依据 AGENTS.md §2 决定二）
%  ⚠️ 本注释的上一版自相矛盾（标题称删2块、实列3个、块号写错），已作废重写。
%
%  【保留 12 块】全部保留：
%   ✓ 块1　{哈密顿--凯莱定理：最小多项式为什么存在
%   ✓ 块2　{最小多项式的定义与四条基本性质
%   ✓ 块3　{最小多项式的计算规则：与准对角、Jordan 块的关系
%   ✓ 块4　{最小多项式无重根 $\iff$ 可相似对角化
%   ✓ 块5　{不变子空间：可对角化的几何解释
%   ✓ 块6　{不变子空间与矩阵形状：分块三角、准对角
%   ✓ 块7　{根子空间分解：一般情形下能保证的只是准对角
%   ✓ 块8　{若尔当标准形的存在与唯一性（北大 §7.8 定理 13、14）
%   ✓ 块9　{相似不变量的完整清单（以及哪一条才够用）
%   ✓ 块10　{实对称矩阵：特征值全实、特征向量正交、必可正交对角化
%   ✓ 块11　{实对称矩阵三条性质的证明要点（每条一个恒等式）
%   ✓ 块12　{求正交矩阵 $T$ 使 $T^{\mathrm{T}}AT$ 成对角形的三步法
%
%  【删除 1 块】
%   ✗ 原块13「正交变换化二次型为标准形——定理 7 的二次型版本」
%      理由：属**第 9 章「二次型」**主题（标准形/规范形/惯性定理），在第 8 章会造成主题越界；
%      第 9 章将由北大 §5.2–§5.4 与樊 §8.2–§8.4 完整覆盖。
%
%  【两处"看似重复实则保留"的说明】
%   · 块1「哈密顿–凯莱定理：最小多项式为什么存在」—— 第 7 章 deep_ch07 已收 H-C 定理，
%     但那一处是把 H-C 当**计算工具**（求 $A^k$、求 $A^{-1}$）；
%     本块用它**论证最小多项式为何存在**，是第 8 章最小多项式理论的入口。
%     子代理曾指出"若删块1，块2 正文并不含存在性论证"——属实，故保留块1。
%   · 块12「求正交矩阵 $T$ 的三步法」—— 对应考纲五.3「掌握实对称矩阵的特征值和特征向量的性质」，
%     是操作层必需，不删。
%
%  【保留重点（对齐考纲五.2/五.3）】
%   最小多项式的定义/性质/计算规则、**无重根 ⟺ 可相似对角化**、
%   不变子空间与可对角化的几何解释、**根子空间分解**、**若尔当标准形的存在与唯一性**、
%   **相似不变量的完整清单**、实对称矩阵三性质及其证明要点、**求正交矩阵 T 的三步法**。
% ===========================================================================
